Вывести на экран ряд натуральных чисел от минимума до максимума с шагом. Например, если минимум $10$, максимум $35$, шаг $5$, то вывод должен быть таким: $10 \ 15 \ 20 \ 25 \ 30 \ 35$. Минимум, максимум и шаг указываются пользователем (считываются с клавиатуры).
Даны два целых числа $A$ и $B$ ($A < B$). Вывести в порядке возрастания все целые числа, расположенные между $A$ и $B$ (включая сами числа $A$ и $B$), а также количество $N$ этих чисел.
Даны два целых числа $A$ и $B$ ($A < B$). Вывести в порядке убывания все квадраты целых чисел, расположенных между $A$ и $B$ (включая числа $A$ и $B$), а также количество $N$ этих чисел.
Дано вещественное число --- цена $1$ кг конфет. Вывести стоимость $1, 2, \dots, 10$ кг конфет.
Дано вещественное число --- цена $1$ кг конфет. Вывести стоимость $0.1, 0.2, \dots, 1$ кг конфет.
Дано вещественное число --- цена $1$ кг конфет. Вывести стоимость $1.2, 1.4, \dots, 2$ кг конфет.
Даны два целых числа $A$ и $B$ ($A < B$). Найти сумму всех целых чисел от $A$ до $B$ включительно.
Даны два целых числа $A$ и $B$ ($A < B$). Найти произведение всех целых чисел от $A$ до $B$ включительно.
Даны два целых числа $A$ и $B$ ($A < B$). Найти сумму квадратов всех целых чисел от $A$ до $B$ включительно.
Дано целое число $N$ ($> 0$). Найти сумму $1 + 1/2 + 1/3 + \dots + 1/N$ (вещественное число).
Дано целое число $N$ ($> 0$). Найти сумму $N^2 + (N + 1)^2 + (N + 2)^2 + \dots + (2 \cdot N)^2$ (целое число).
Дано целое число $N$ ($> 0$). Найти произведение $1.1 \cdot 1.2 \cdot 1.3 \cdot \dots$ ($N$ сомножителей).
Дано целое число $N$ ($> 0$). Найти значение выражения $1.1 - 1.2 + 1.3 - \dots$ ($N$ слагаемых, знаки чередуются). Условный оператор не использовать.
Дано целое число $N$ ($> 0$). Найти квадрат данного числа, используя для его вычисления следующую формулу: $N^2 = 1 + 3 + 5 + \dots + (2 \cdot N - 1)$. После добавления к сумме каждого слагаемого выводить текущее значение суммы (в результате будут выведены квадраты всех целых чисел от $1$ до $N$).
Дано вещественное число $A$ и целое число $N$ ($> 0$). Найти $A$ в степени $N$: $A^N = A \cdot A \cdot \ldots \cdot A$ (числа $A$ перемножаются $N$ раз).
Дано вещественное число $A$ и целое число $N$ ($> 0$). Используя один цикл, вывести все целые степени числа $A$ от $1$ до $N$.
Дано вещественное число $A$ и целое число $N$ ($> 0$). Используя один цикл, найти сумму $1 + A + A^2 + A^3 + \dots + A^N$.
Дано вещественное число $A$ и целое число $N$ ($> 0$). Используя один цикл, найти значение выражения $1 - A + A^2 - A^3 + \dots + (-1)^N \cdot A^N$. Условный оператор не использовать.
Дано целое число $N$ ($> 0$). Найти произведение $N! = 1 \cdot 2 \cdot \ldots \cdot N$ ($N$--факториал). Чтобы избежать целочисленного переполнения, вычислять это произведение с помощью вещественной переменной и вывести его как вещественное число.
Дано целое число $N$ ($> 0$). Используя один цикл, найти сумму $1! + 2! + 3! + \dots + N!$ (выражение $N!$ --- $N$--факториал --- обозначает произведение всех целых чисел от $1$ до $N$: $N! = 1 \cdot 2 \cdot \ldots \cdot N$). Чтобы избежать целочисленного переполнения, проводить вычисления с помощью вещественных переменных и вывести результат как вещественное число.
Дано целое число $N$ ($> 0$). Используя один цикл, найти сумму $1 + 1/(1!) + 1/(2!) + 1/(3!) + \dots + 1/(N!)$ (выражение $N!$ --- $N$--факториал --- обозначает произведение всех целых чисел от $1$ до $N$: $N! = 1 \cdot 2 \cdot \ldots \cdot N$). Полученное число является приближенным значением константы $e = \exp(1)$.
Дано вещественное число $X$ и целое число $N$ ($> 0$). Найти значение выражения $1 + X + X^2/(2!) + \dots + X^N/(N!)$ ($N! = 1 \cdot 2 \cdot \ldots \cdot N$). Полученное число является приближенным значением функции $\exp$ в точке $X$.
Дано вещественное число $X$ и целое число $N$ ($> 0$). Найти значение выражения $X - X^3/(3!) + X^5/(5!) - \dots + (-1)^N \cdot X^{2 \cdot N + 1}/((2 \cdot N + 1)!)$ ($N! = 1 \cdot 2 \cdot \ldots \cdot N$). Полученное число является приближенным значением функции $\sin$ в точке $X$.
Дано вещественное число $X$ и целое число $N$ ($> 0$). Найти значение выражения $1 - X^2/(2!) + X^4/(4!) - \dots + (-1)^N \cdot X^{2 \cdot N}/((2 \cdot N)!)$ ($N! = 1 \cdot 2 \cdot \ldots \cdot N$). Полученное число является приближенным значением функции $\cos$ в точке $X$.
Дано вещественное число $X$ ($|X| < 1$) и целое число $N$ ($> 0$). Найти значение выражения $X - X^2/2 + X^3/3 - \dots + (-1)^{N - 1} \cdot X^N/N$. Полученное число является приближенным значением функции $\ln$ в точке $1 + X$.
Напишите программу, доказывающую или проверяющую, что для множества натуральных чисел выполняется равенство: $1 + 2 + \dots + n = n(n + 1)/2$, где $n$ --- любое натуральное число.
Найти сумму $n$ элементов следующего ряда чисел: $1, -0.5, 0.25, -0.125, \dots$ Количество элементов ($n$) вводится с клавиатуры.
Вывести на экран столько элементов ряда Фибоначчи, сколько указал пользователь. Например, если на ввод поступило число $6$, то вывод должен содержать шесть первых чисел ряда Фибоначчи. Ряд Фибоначчи --- это последовательность натуральных чисел, где каждое последующее число является суммой двух предыдущих: $1 \ 1 \ 2 \ 3 \ 5 \ 8 \ 13 \ 21 \ 34 \ 55 \ 89 \dots$
Вычислить: $(1 + 2) \cdot (1 + 2 + 3) \cdot \dots \cdot (1 + 2 + \dots + 10)$.
Вычислить: $1 + 2 + 4 + 8 + \dots + 2^{10}$.
